\documentclass[10pt,a4paper]{article}
\usepackage{amsmath,amssymb,amsthm,graphicx}

\newtheorem{theorem}{Theorem}[section]

\newtheorem{lemma}[theorem]{Lemma}
\newtheorem{proposition}[theorem]{Proposition}
\newtheorem{observation}[theorem]{Observation}
\newtheorem{fact}[theorem]{Fact}
\newtheorem{corollary}[theorem]{Corollary}

\theoremstyle{definition}
\newtheorem{definition}[theorem]{Definition}
\newtheorem{conjecture}[theorem]{Conjecture}
\newtheorem{example}[theorem]{Example}
\newtheorem{remark}[theorem]{Remark}
\newtheorem{question}[theorem]{Open Question}

\DeclareMathOperator{\GL}{GL}
\DeclareMathOperator{\Span}{Span}
\DeclareMathOperator{\Id}{Id}
\DeclareMathOperator{\Hom}{Hom}
\DeclareMathOperator{\rk}{rk}
\newcommand{\F}{\mathbb{F}}
\newcommand{\R}{\mathbb{R}}
\newcommand{\E}{\mathbb{E}}
\newcommand{\diag}{\text{diag}}

\title{On the algebraic proof complexity of Tensor Isomorphism}
\usepackage{hyperref}
\usepackage{arydshln}
\title{Polynomial-calculus degree lower bounds for the rank and inversion principles}
\author{Nicola Galesi, Joshua A. Grochow, Toniann Pitassi, Adrian She}
\date{doi:10.4230/LIPIcs.CCC.2023.4; CC BY 4.0}
\begin{document}
\maketitle
\noindent\textit{Editorial note.} Public sources contain only the corresponding official published complete author material; native TeX served as editable transcription aid. Raw native source/tar remains private. Original theorem numbering preserved with section counters. Rank and inversion consequences counted as a single result, and the alternative Random3XOR tensor lower-bound proof is not a new record.\par\medskip
\noindent\textit{Editorial note.} References to unselected original sections are printed with their accurate original article chapter numbers; selected section headers and original labels remain in place. No absent section is presented as an included proof.\par\medskip
\noindent\textit{Allowed standard prerequisites (editorial).}\begin{itemize}
\item Known graph-isomorphism and functional-pigeonhole polynomial-calculus lower bounds; effective Nullstellensatz and tensor-rank inheritance. Exact original cited statements retained where used; standard earlier literature not recursively reproved
\end{itemize}\medskip


\begin{abstract}
The \textsc{Tensor Isomorphism} problem (\textsc{TI}) has recently emerged as having connections to multiple areas of research within complexity and beyond, but the current best upper bound is essentially the brute force algorithm. Being an algebraic problem, \textsc{TI} (or rather, proving that two tensors are \emph{non}-isomorphic) lends itself very naturally to algebraic and semi-algebraic proof systems, such as the Polynomial Calculus (PC) and Sum of Squares (SoS). For its combinatorial cousin \textsc{Graph Isomorphism}, essentially optimal lower bounds are known for approaches based on PC and SoS (Berkholz \& Grohe, SODA '17). Our main results are an $\Omega(n)$ lower bound on PC degree or SoS degree for \textsc{Tensor Isomorphism}, and a nontrivial upper bound for testing isomorphism of tensors of bounded rank.

We also show that PC cannot perform basic linear algebra in sub-linear degree, such as comparing the rank of two matrices, or deriving $BA=I$ from $AB=I$. As linear algebra is a key tool for understanding tensors, we introduce a strictly stronger proof system, PC+Inv, which allows as derivation rules all substitution instances of the implication $AB=I \rightarrow BA=I$. We conjecture that even PC+Inv cannot solve  \textsc{TI} in polynomial time either, but leave open getting lower bounds on PC+Inv for any system of equations, let alone those for \textsc{TI}. We also highlight many other open questions about proof complexity approaches to \textsc{TI}.
\end{abstract}



\section{Introduction}
Tensors have rapidly emerged as a fundamental data structure and key mathematical object of the 21st century. They play key roles in many different areas of science, engineering, and mathematics, from quantum mechanics and general relativity to neural networks \cite{TNN} and mechanical engineering. They arise in theoretical computer science in many ways, including from (post-quantum) cryptography \cite{Pat96, JQSY19}, derandomization, \textsc{Matrix Multiplication}, \textsc{Graph Isomorphism} \cite{GQ}, and several different parts of Geometric Complexity Theory. 

The fundamental notion of equivalence between tensors is that of isomorphism: two tensors are isomorphic if one can be transformed into the other by an invertible linear change of basis in each of the corresponding vector spaces. For example, two 2-tensors (=matrices) $M,M'$ are equivalent under this notion if there are invertible matrices $X,Y$ such that $XMY = M'$; similarly, two 3-tensors, represented by 3-way arrays $T_{ijk}, T'_{ijk}$ are isomorphic if there are three invertible matrices $X,Y,Z$ such that 
\begin{equation} \label{eq:TI}
\sum_{ijk} X_{ii'} Y_{jj'} Z_{kk'} T_{ijk} = T'_{i'j'k'}
\end{equation}
for all $i',j',k'$. The problem of (3-)\textsc{Tensor Isomorphism} (\textsc{TI}) is: given two such 3-way arrays, to decide if they are isomorphic. 

Over finite fields, two different versions of \textsc{TI} sandwich the complexity of its more famous cousin, \textsc{Graph Isomorphism}. Namely, as presented above, \textsc{GI} reduces to \textsc{TI}. In the other direction, over a finite field $\F=\F_{p^a}$, one can take an $n \times n \times n$ tensor and list it out ``verbosely'', as a set of $p^{an}$ many $n \times n$ matrices over $\F$; the isomorphism problem for such verbosely given tensors is equivalent to \textsc{Group Isomorphism} for a certain class of $p$-groups, widely believed to be the hardest cases of \textsc{Group Isomorphism} in general. As such, this verbose version of \textsc{TI} reduces to \textsc{GI}. Furthermore, with Babai's quasi-polynomial-time algorithm \cite{Babai}, the running times are quite close: $N^{O(\log N)}$ for \textsc{VerboseTI} and $N^{O(\log^2 N)}$ for \textsc{GI} (the exponent of the exponent was worked out by Helfgott \cite{helfgott}). Thus \textsc{TI} stands as a key obstacle to putting \textsc{GI} into $\mathsf{P}$.


In this paper, we initiate the study of (algebraic) proof complexity approaches to proving that two tensors are non-isomorphic.
Lower bounds on the Polynomial Calculus proof system imply lower bounds on Gr\"{o}bner basis techniques, and the latter are some of the leading methods for solving $\mathsf{TI}$-complete problems in cryptanalysis, e.g., \cite{DBLP:conf/eurocrypt/TangDJPQS22, DBLP:conf/eurocrypt/FaugereP06}.
In the context of \textsc{GI}, proof complexity plays an important role, through its connection with the Weisfeiler--Leman (WL) algorithm. Although this algorithm does not, on its own, solve \textsc{GI} in polynomial time \cite{CFI}, it is a key subroutine in many of the best algorithms for \textsc{GI}, both in theory \cite{Babai} and in practice (see \cite{mckay, MP}). And the picture that has emerged is that some proof systems for \textsc{GI} are known to be equivalent in power to WL \cite{AM, BerkholzG15}, and some lower bounds on proof systems are closely related to lower bounds for WL \cite{SSC, DBLP:conf/soda/ODonnellWWZ14, BerkholzG15}. Versions of WL for groups, and in particular finite $p$-groups---and hence, by the connection above, tensors over finite fields---have only recently begun to be explored \cite{BGLQW, BS1, BS2, CL}.

\subsection{Main results}

We focus on the Polynomial Calculus (PC, or Gr\"{o}bner) proof system \cite{CEI}, though our results will also hold for semi-algebraic proof systems such as Sum-of-Squares \cite{lasserre} as well. PC is used to show that a system of polynomial equations over a field $\F$ is unsatisfiable over the algebraic closure $\overline{\F}$, by deriving from the system of equations, in a line-by-line fashion, the contradiction $1=0$. The degree of a PC proof is the maximum degree of any line appearing in the proof, and it is a fundamental result that PC proofs of constant degree can be found in polynomial time \cite{CEI}. Much as WL informally ``captures all combinatorial approaches'' to \textsc{GI}, PC informally ``captures all approaches based on Gr\"{o}bner bases'' to showing that a system of polynomial equations is unsatisfiable. 

The systems of equations we study are, for two non-isomorphic tensors $T,T'$, the equations (\ref{eq:TI}) along with new matrices $X',Y',Z'$, and equations saying that these are the inverses of $X,Y,Z$, resp., viz.: $XX'=X'X = \Id$, and similarly for the others. The reason for introducing these new matrices, despite their not appearing in (\ref{eq:TI}), is that these invertibility equations are only degree 2. In contrast, if we instead used the determinant to indicate that $X$ was invertible, then our starting equations would have degree $n+1$, rather than constant degree $\leq 3$. Since the main complexity measure we study on PC is degree, having starting equations of degree $n$ would make it difficult to make meaningful lower bound statements.

Our first main result is (two proofs of) a lower bound on such techniques.

\begin{theorem} \label{thm:lowerMain}
Over any field, there are instances of $n \times n \times n$ \textsc{Tensor Isomorphism} that require PC degree $\Omega(n)$ to refute. Over $\R$, they also require Sum-of-Squares degree $\Omega(n)$ to refute.
\end{theorem}

The preceding goes by reduction from known lower bounds on PC for \textsc{Graph Isomorphism} \cite{BerkholzG15, BerkholzG16}, but has the disadvantage (from the tensor point of view) that the resulting tensors are quite sparse: in one direction, one of the slices is supported on an $\Omega(n) \times n$ matrix and all the others slices have support size 1. In a second proof (Section~6 (in the original article)), we get a polynomially worse lower bound $\Omega(\sqrt[4]{n})$, but with a reduction from \textsc{Random 3XOR} that is more direct. Indeed, we show that 3XOR itself can be viewed as a particular instance of a tensor problem \emph{without} gadgets; gadgets are only then needed to reduce from that tensor problem to \textsc{Tensor Isomorphism} itself. In contrast, the lower bounds on PC for \textsc{GI} (\emph{ibid.}) already use the Cai--F\"{u}rer--Immerman gadgets \cite{CFI} to reduce from XOR-SAT, and then even further gadgets are needed to reduce from \textsc{GI} to \textsc{TI}. 

Our technical contributions in the above theorem are thus three-fold: 
\begin{enumerate}
\item We show that the known reductions from \textsc{GI} to \textsc{TI} can be carried out in low-degree PC; 
\item We realize 3XOR very naturally as a tensor problem; and 
\item We give new reductions from 3XOR, through a series of tensor-related problems, to \textsc{TI}, that work as many-one reductions of the decision problems that can be carried out in low-degree PC.
\end{enumerate}

Complementing our lower bound, we also show that tensors of low rank are comparatively easy to test for (non)-isomorphism. Here, one of our upper bounds is in the weaker Nullstellensatz proof system (giving a stronger upper bound than only a PC upper bound). In the Nullstellensatz proof system, a proof that a system of equations $f_1 = \dotsb = f_m = 0$ is unsatisfiable consists of polynomials $g_i$ such that $\sum g_i f_i = 1$, and the Nullstellensatz degree is the maximum degree of any $g_i f_i$. The PC degree is always at most the Nullstellensatz degree, and the gap between the two can be nearly maximal for Boolean equations: $O(1)$ versus $\Omega(n / \log n)$ \cite{BCIP}. (For Boolean equations, there is always an $O(n)$ upper bound, though this does not apply to \textsc{TI}, see Remark~\ref{rmk:boolean} below).

\begin{theorem} \label{thm:upper}
Over any field, the Nullstellensatz degree of refuting isomorphism of two $n \times n \times n$ tensors of tensor rank $\leq r$ is at most $2^{O(r^2)}$. If working over a finite field $\F_q$ and including the equations $x^q - x$, the PC degree is at most $O(qr^2)$. 

In particular, isomorphism of constant-rank tensors can be decided in polynomial time.
\end{theorem}

\begin{remark} \label{rmk:boolean}
In many settings in proof complexity, Boolean axioms such as $x_i^2 = x_i$ or $x_i^2 = 1$ are included among the system of equations, and all such unsatisfiable systems of equations can be refuted in degree $O(n)$ ($n=$\# variables). If this were the case here, the above would only be interesting for very small values of $r$. In contrast, the equations for \textsc{TI} do not include any such Boolean axioms, and as such the naive degree upper bound is exponential in the number of variables. For $n \times n \times n$ tensors, this gives an upper bound of $2^{O(n^2)}$ \cite{sombra}, and thus, Theorem~\ref{thm:upper} gives nontrivial upper bounds all the way up to $r \leq n$. (We note that $n \times n \times n$ tensors can have rank up to $\Theta(n^2)$ \cite{lickteigTypical}.) The proof of Theorem~\ref{thm:upper} shows that for rank-$r$ tensors, \textsc{TI} can essentially be reduced to a system of equations in only $O(r^2)$ variables.
\end{remark}

\begin{remark} \label{rmk:rank}
For fixed $r$, testing if an $n \times n \times n$ tensor has rank $\leq r$ can be done in polynomial time, as follows. This will show that the algorithm of Theorem~\ref{thm:upper} genuinely solves the decision problem, and not just a promise problem. Given an $n \times n \times n$ tensor $T$, consider its three $n \times n^2$ flattenings. Use Gaussian elimination to put each such flattening, separately, into reduced row echelon form. If any of these flattenings has rank $> r$, reject. Otherwise, we get from this a list of $3r$ vectors $u_1, \dotsc, u_r, v_1, \dotsc, v_r, w_1, \dotsc, w_r$, such that $T$ lives in the $r \times r \times r$-dimensional space $\Span\{u_1, \dotsc, u_r\} \otimes \Span\{v_1, \dotsc, v_r\} \otimes \Span\{w_1, \dotsc, w_r\}$. Now in this space we can write down the Brent equations \cite{brent} for $T$ to have rank $\leq r$, which will be $r^3$ cubic equations in $3r^2$ variables (Brent's equations \cite[(5.06)]{brent} were specifically for the matrix multiplication tensor, but analogous equations are easily constructed for arbitrary tensors using the same idea). Since $r$ is constant, these equations may be solved in polynomial time (here we assume that we are either working over a finite field, a finite-degree extension of the rationals---see, for example, Grigoriev \cite{grigoriev}---or in the BSS model over an arbitrary field).
\end{remark}

Lastly, one may wonder why we focus on 3-\textsc{Tensor Isomorphism}, and not some of its many related variants. Indeed, just as there are other equivalence notions for matrices---such as conjugacy $XMX^{-1}$ and congruence $XMX^T$---there are many different kinds of multilinear objects that can be represented by multi-way arrays, including tensors, homogeneous polynomials (commutative or noncommutative), alternating matrix spaces, multilinear maps, and so on, each with their own corresponding notion of isomorphism. While these problems are indeed distinct, they are all equivalent under polynomial-time isomorphisms \cite{FGS, GQ}; such problems are called $\mathsf{TI}$-complete. Even isomorphism of $k$-way tensors (for any fixed $k \geq 3$) is equivalent to isomorphism of 3-tensors \cite{GQ}. This partially justifies our focus on \textsc{3-Tensor Isomorphism}. In the course of proving our reductions for the results stated above, we use many of the gadgets from \cite{FGS, GQ}, and show that such uses also often yield proof complexity reductions as well. Because of the variety of gadgets used in our reductions, we believe that many, if not all, of the gadgets from those results would also yield proof complexity reductions, so the proof complexity of all the known $\mathsf{TI}$-complete problems should be polynomially related.

\subsection{Comparison with linear algebra, a new proof system, and a conjecture} \label{sec:conj}
As linear algebra is part of the core toolkit for understanding tensors, it is natural to wonder how linear algebra can help in algebraic proof complexity approaches to \textsc{TI}. We believe that even if it had the ``full power'' of linear algebra at its disposal ``for free,'' PC could still not solve \textsc{TI} efficiently. We begin to make this precise in this section.

Some basic derivations in linear algebra are to relate the ranks of two matrices and to derive $BA=I$ from $AB=I$ (the Inversion Principle, one of the so-called ``hard matrix identities'' \cite{SC}, only recently shown to have short $\mathsf{NC}^2$-Frege proofs \cite{HT}). Soltys \cite{soltys} and Soltys \& Cook \cite{SC} discuss the relationship between these and other standard implications in linear algebra. We show that PC is not strong enough to prove these in low-degree:


\begin{theorem}
The unsatisfiable system of equations $XY=\Id_n$ where $X$ is $n \times r$ and $Y$ is $r \times n$ with $1 \leq r < n$, requires degree $\geq r/2+1$ to refute in PC, over any field.
\end{theorem}
We refer to this system of equations as the Rank Principle, as refuting them amounts to showing that $\rk \Id_n > r$. 

\begin{theorem}
Any PC derivation of $BA = I$ from $AB=I$, where $A,B$ are $n \times n$ matrices with $\{0,1\}$ entries, requires degree $\geq n/2+1$, over any field.
\end{theorem}


We also observe that the Rank Principle can be derived in low degree from the Inversion Principle. 

Although it remains open whether the Inversion Principle is ``complete'' for linear-algebraic reasoning (see \cite{soltys, SC}), we introduce the proof system PC+Inv in an attempt to capture some linear-algebraic reasoning that seems potentially useful for \textsc{TI}. PC+Inv has all the same derivation rules as PC, but in addition, for any square matrices $A,B$ (whose entries may themselves be polynomials---that is, we allow substitution instances), we have the rule
\[
\frac{AB=I}{BA=I}.
\]
where the antecedent represents the set of $n^2$ equations corresponding to $AB=I$, and similarly the consequent denotes the set of $n^2$ equations $BA=I$ (see \ref{sec:prelim:inv} for more details). Degree is still measured in the usual way, but this rule lets us ``cut out'' the high-degree proof that would usually be required to derive $BA=I$ from $AB=I$. We now formalize our intuition that linear algebra should not suffice to solve \textsc{TI} efficiently in the following:

\begin{conjecture} \label{conj:inv}
\textsc{Tensor Isomorphism} for $n \times n \times n$ tensors requires degree $\Omega(n)$ in PC+Inv, over any field.
\end{conjecture}

Despite the conjecture, we do not yet know how to prove lower bounds on PC+Inv for any unsatisfiable system of equations, let alone those coming from \textsc{TI}. Mod $p$ counting principles (for $p$ different from the characteristic of the field) strike us as potentially interesting instances to examine for PC+Inv lower bounds, before tackling a harder problem like \textsc{TI}. In the final section, we highlight many other open questions around the proof complexity of \textsc{TI}.

\subsection{Organization}
In Section~\ref{sec:prelim} we cover preliminaries. In Section~\ref{sec:matrix} we prove the lower bounds on linear algebraic principles just discussed. In Section~4 (in the original article) we prove the upper bound for isomorphism of bounded rank tensors (Theorem~\ref{thm:upper}). In Section~5 (in the original article) we prove Theorem~\ref{thm:lowerMain} by reduction from \textsc{GI}. In Section~6 (in the original article) we prove the polynomially related lower bound by direct reduction from \textsc{Random 3XOR}.



\section{Preliminaries} \label{sec:prelim}


\subsection{Proof systems}
All our rings are commutative and unital. {\em Polynomial calculus} (PC) is a proof system to prove that a given system of (multivariate) polynomial equations ${\cal P}$ over a field $\mathbb F$ of the form $p=0$, has no solution over the algebraic closure $\overline{\F}$ (i.e. the system is unsolvable). We usually shorten the polynomial equation $p=0$ to just $p$.
The derivation rules of the system are the following one:
$$
\frac{p}{xp} \; \mbox{(multiplication)}, \quad \frac{p\quad q}{ap+bq} \; \mbox{(linear combination)} 
$$
where $x$ is any formal variable, $a,b \in \mathbb{F}$ and $p,q$ are polynomials over $\mathbb F$.  

When refuting Boolean systems of equations it is common to include the Boolean axioms $x_i^2 - x_i$. Because we do \emph{not} always include these (esp. for \textsc{TI}) we are explicit about our use of these, but do not assume they are built into the proof system---that is, if we are assuming them as axioms, we say so.

A {\em PC derivation (or proof)}  of a  polynomial $q$ from a set of polynomials ${\cal P}$ is a sequence of polynomial equations $p_1,\ldots, p_{m}$ ending with the polynomial $q$ (so $p_m$ is $q$) and where each $p_i$, $i \in [m]$, is either an {\em axiom} $p$ for $p \in  {\cal P}$, or is obtained from previous equations in the refutation by multiplication or linear combination.  We denote this by writing  ${\cal P}\vdash q$.  Observe that if $p$  is derivable in  PC and $q$ is a polynomial then, by repeated applications of multiplication and linear combination rules, we can derive
$pq$. We often use this generalization of the multiplication in our proofs without mention.  

A {\em PC refutation}  is just a PC proof  of the polynomial $1$.   The {\em degree} of a PC derivation is the maximal degree of a polynomial used in the proof. The {\em size} of a polynomial $p$ is the number of  terms in $p$.  The {\em size} a PC derivation $p_1,\ldots, p_m$  is the sum of the sizes of the polynomials $p_1,\ldots,  p_m$. 


\medskip

For our upper bound in Theorem~\ref{thm:upper}, we also consider another algebraic proof system, known as {\em Nullstellensatz} (NS), to certify unsolvability of sets of polynomial equations. Nullstellensatz  is defined in a static form as follows:  a  refutation of a list  ${\cal P}=(p_1,\ldots,p_m)$ of polynomial equations over variables $x_1\ldots,x_n$ 
is given by the list of polynomials  ${\cal Q}=(q_1,\ldots q_m)$  such that
$$
\sum_{i \in [m]} p_iq_i =1
$$

The {\em degree} of a NS refutation is the maximal degree of a polynomial in ${\cal P} \cup {\cal Q}$. The size of $NS$ proof is the sum of the number of monomials appearing in the polynomials
$q_1,\ldots, q_m$.

\medskip

{\em Sum-of-Squares} (SOS) is a static proof system for certifying the unsolvability of systems of polynomial equations and polynomial inequalities, where polynomials are usually over the ring $\mathbb R[x_1\ldots, x_n]$.    

A polynomial $p$ is a {\em sum-of-squares} polynomial if it is in the form $p=\sum_{i} r_i^2$ and the $r_i$'s  are polynomials as well.  Given a system made by a set of polynomial equations  
 ${\cal P} =\{p_1=0,\ldots p_m=0\}$ and a set ${\cal Q}=\{q_1\geq 0,\ldots q_k\geq 0\}$ of polynomial inequalities, a {\em sum-of-squares} proof of the polynomial inequality $p\geq 0 $ from ${\cal P} \cup {\cal Q}$  is  given by the  formal identity
 $$
 p= s_0+\sum_{i\in [k]}s_iq_i + \sum_{j \in [m]}t_jp_j
 $$
where $s_0,s_1,\ldots, s_k$ are sum-of-squares polynomials, while $t_1,\ldots,t_m$ are arbitrary polynomials.  When the system  ${\cal P} \cup {\cal Q}$  is unsatisfiable, a {\em refutation} of ${\cal P} \cup {\cal Q}$ is a proof of the inequality $-1\geq 0$, that is for  $p$ the constant polynomial $-1$. The {\em degree} of the proof is 
the $\max \{\deg(p),\deg(s_0), \deg(s_i)+\deg(q_i), \deg(t_j)+\deg(p_j) | i \in [k], j \in[m]\}$.   


\medskip 

\begin{definition}[PC reduction between systems of polynomials, cf. {\cite[Sec.~3]{DBLP:journals/jcss/BussGIP01}}] \label{def:reduction}
Let $P(x_1,\ldots,x_n)$ and $Q(y_1,\ldots, y_m)$ be two sets of polynomials over a field $\mathbb F$.  
$P$ is \emph{$(d_1,d_2)$-reducible} to $Q$ if:
\begin{enumerate}
\item For each $i \in [m]$ there is a polynomial $r_i(\mathbf{x})$ of degree at most $d_1$  (which we think of as defining $y_i$ in terms of the $\mathbf x$ variables);
\item There exists a degree $d_2$ PC derivation of $Q(r_1(\mathbf x),\ldots, r_m(\mathbf x) )$ from polynomials $P(\mathbf x)$.
\end{enumerate}
\end{definition}

\begin{lemma}[{\cite[Lem.~1]{DBLP:journals/jcss/BussGIP01}}]
\label{lem:PCred}
If $P(\mathbf{x})$ is $(d_1, d_2)$-reducible to $Q(\mathbf{y})$ and there is
a degree $d$ PC refutation of $Q(\mathbf{y})$, then there is a degree $\max(d_2, d_1 d)$ 
refutation of $P(\mathbf x)$.
\end{lemma}

In their paper, they typically only applied this to systems of equations which were known to be unsatisfiable (such as PHP and Tseitin tautologies), whereas in our paper we have several situations we want to combine the above notion together with the usual notion of many-one reduction. We encapsulate this in the following definition. We say a decision problem $\Pi$ is a \emph{polynomial solvability problem} over a field $\F$ if all valid instances of the problem are systems of polynomial equations over $\F$, and the problem is to decide whether such a system of equations has solutions over the algebraic closure $\overline{\F}$. Thus, the difference between multiple polynomial solvability problems is just \emph{which} systems of equations are valid inputs.

\begin{definition}[PC many-one reduction] \label{def:reduction2}
Let $\Pi_1, \Pi_2$ be two polynomial solvability problems over a field $\F$. We say that $\Pi_1$ \emph{$(d_1,d_2)$-many-one reduces} to $\Pi_2$ if there is a polynomial-time many-one reduction $\rho$ from $\Pi_1$ to $\Pi_2$, such that for all unsatisfiable instances $\mathcal{F}$ of $\Pi_1$, $\mathcal{F}$ $(d_1, d_2)$-reduces to $\rho(\mathcal{F})$. When this occurs with $d_1,d_2 = O(1)$, we write \[\Pi_1 \leq_m^{PC} \Pi_2.\]
\end{definition}



\subsection{Linear algebra and tensors}
Given three vector spaces $U,V,W$ over a field $\F$, a 3-tensor is an element of the vector space $U \otimes V \otimes W$, whose dimension is $(\dim U)(\dim V)(\dim W)$. If $e_i$ is the $i$-th standard basis vector, then a basis for $U \otimes V \otimes W$ is given by the vectors $\{e_i \otimes e_j \otimes e_k\}$. One may also interpret the symbol $\otimes$ more concretely as the Kronecker product, in which $e_i \otimes e_j \otimes e_k$ represents a 3-way array whose only nonzero entry is in the $(i,j,k)$ position. The vector space of such 3-way arrays (with coordinate-wise addition) is isomorphic to $U \otimes V \otimes W$.

The rank of a tensor $T \in U \otimes V \otimes W$ is the minimum $r$ such that $T = \sum_{i=1}^r u_i \otimes v_i \otimes w_i$ for some vectors $u_i, v_i, w_i$. 

Two $n \times m \times p$ 3-tensors $T,T' \in U \otimes V \otimes W$ are isomorphic if there exist matrices $X \in \GL(U), Y \in \GL(V), Z \in \GL(W)$ such that $(X,Y,Z) \cdot T = T'$, where the latter is shorthand for (\ref{eq:TI}).
If we treat $T,T'$ as given non-isomorphic tensors, then we may treat (\ref{eq:TI}) as a system of equations in the $n^2 + m^2 + p^2$ variables $X_{ii'}, Y_{jj'}, Z_{kk'}$. To enforce that these variable matrices are invertible, we furthermore introduce three additional sets of variables $X', Y', Z'$ meant to be the inverse matrices, and include also the equations 
\[
XX' = X'X = I_n \qquad YY' = Y'Y = I_m \qquad ZZ' = Z' Z = I_p,
\]
where $I_n$ denotes the $n \times n$ identity matrix, which is $\Id_U$ in any basis. (We could have instead introduced new variables such as $\delta$ and the equation $\det(X) \delta = 1$, however, the latter equation is degree $n$, whereas the above equations all have degree $O(1)$, which is more desirable from the point of view of algebraic proof complexity.)

\subsection{Polynomial encodings and the inversion principle} \label{sec:prelim:inv}
Some  principles of linear algebra can be formulated as tautologies in propositional logic and therefore also as a set of polynomial equations.
In this paper we preliminarily consider two such principles. 

\medskip

\noindent {\bf Rank Principle.} As a first example we consider a set of unsatisfiable polynomials encoding the principle that  the product of a $n\times r$ matrix $X$ by a $r\times n$ matrix $Y$ cannot be the identity matrix whenever $r<n$.  We consider variables $x_{i,k},y_{j,k}$ for $i,j\in [n]$ and $k\in [r]$, where $r< n$ to encode $X$ and $Y$. Then the polynomial encoding is:
$$\mathbb I(r,n) := \sum_{k\in[r]} x_{i,k}y_{j,k}-\delta_{i,j} \quad \quad i,j\in [n]$$
where $\delta_{i,j}=1$ if $i=j$ and $0$ otherwise. This set of polynomials is clearly unsatisfiable as long as $r<n$. 

\medskip

\noindent {\bf Inversion Principle.} The second principle encodes the invertibility of a square $n\times n$ matrix $A$,  expressing the tautology that
 $AB=I \rightarrow BA=I$ where $A,B$ are $n\times n$ matrices and $I$ is the identity matrix. Stephen A. Cook suggested this principle as a tautology that may be hard to prove in several proof systems. 
 
Let $a_{i,j},b_{i,j}$ be formal variables 
 encoding respectively  the $(i,j)$-th entries of $A$ and $B$.  We represent the fact that $AB=I$ as the set of degree $2$ polynomials

$$
\sum_{k \in [n]} a_{i,k}b_{k,j} - \delta_{i,j}\qquad i,j \in [n],
$$
where $\delta_{i,j}=1$ if $i=j$ and $0$ otherwise. We denote this set of polynomials by $AB=I$. In Section \ref{sec:matrix}, we study the degree complexity of  $AB=I \vdash BA=I$, that is of PC derivations of the polynomials $BA=I$ from the polynomials $AB=I$. 

\medskip

In view of the results we obtain in Section \ref{sec:matrix}, in Section~\ref{sec:conj} we considered a  {\em polynomial rule schema} of the form 
$$
\frac{AB=I}{BA=I}
$$
which we call the {\em Inversion Rule} (INV) meant to be added to $PC$ as an extra rule. We make this slightly more precise here.

A polynomial instantiation $\tau$ of the polynomials $AB=I$ is a substitution of polynomials $p_{i,j},q_{i,j}$ to variables $a_{i,j}$ and $b_{i,j}$.   
In PC+INV a polynomial $p$ is derivable from a set of polynomials
${\cal P}$ if 
\begin{enumerate}
\item $p$ is an axiom, or $p \in {\cal P}$;
\item $p$ is obtained by multiplication or linear combination from previous polynomials in the proof;
\item $p$ is a polynomial among a polynomial instantiation $\tau$ of $BA=I$, given that among the polynomials previously derived  in the proof there are all the polynomials 
forming the instantiation $\tau$ of $AB=I$.
\end{enumerate}

\medskip

\noindent {\bf Pigeonhole Principle.} An important role in proving the results in Section \ref{sec:matrix} is played by the well-known {\em Pigeonhole principle} stating that  any function $f$ from $[n]$ to $[r]$ with $r<n$ has a collision, that is there are $i \not= i' \in [n]$ and a $j\in[r]$ such that $f(i)=f(i')=j$. $PHP^n_r$ is the set of polynomials:

$$
\sum_{k\in [r]}p_{i,k} -1, \mbox{ for } i \in [n], \qquad \qquad p_{i,k}p_{j,k}, \mbox{ for } i \not= j \in [n], k \in [r] \qquad p_{ij}^2 - p_{ij} \text{, for $i \in [n], j \in [r]$} 
$$

Razborov \cite{DBLP:journals/cc/Razborov98} additionally included the ``functional equations'' (encoding that each pigeon cannot be matched to more than one hole):
\[
p_{i,k} p_{i,k'}, \mbox{ for } i \in [n], k \neq k' \in [r].
\]

\medskip






\section{Linear algebra warm-up: PC for matrices}
\label{sec:matrix}
Two matrices $M,M' \in U \otimes V$ are isomorphic as tensors if they are equivalent as matrices, meaning under left- and right-multiplication by invertible matrices $X \in \GL(U), Y \in GL(V)$, that is, 
\[
XMY = M'. 
\]

Since we want $X,Y$ to be invertible, we also introduce variable matrices $X', Y'$ as before, together with the equations
\[
XX' = X'X = \Id_U \qquad YY' = Y'Y = \Id_V.
\]
Then by left multiplying our initial matrix equation by $Y'$, we may replace it with the new matrix equation
\[
XM = M'Y'.
\]
The latter has the advantage of being linear in $X$ and $Y'$, but the quadratic equations $XX' =  \Id_U, YY' =  \Id_V$ still make even this case not totally obvious.

\subsection{A trick for PC degree} \label{sec:trick}
If our focus is on PC \emph{degree}, we note that the degree of the equations is unchanged if we first left- or right-multiply $M,M'$ by invertible scalar matrices. For example, if we replace $M$ by $\overline{M} = AMB$ with $A,B \in \GL(U)$, then we may replace $X$ by $\overline{X} := XA^{-1}$, $Y$ by $\overline{Y} := B^{-1} Y$. Then we have $\overline{M} \cong M$, so $\overline{M} \cong M'$ iff $M \cong M'$. Furthermore, since the transformation $X \mapsto XA^{-1}$, $Y \mapsto B^{-1} Y$ is linear and invertible, any PC proof that $\overline{M} \not\cong M'$ can be transformed by the inverse linear transformation into a PC proof that $M \not\cong M'$ of the same degree.

Now, for matrices under this equivalence relation, we have a normal form, namely every matrix $M$ is equivalent to a diagonal matrix with $\rk(M)$ 1s on the diagonal and all the remaining entries $0$, that is, $\sum_{i=1}^{\rk(M)} e_i \otimes e_i = I_r \oplus 0$, where the latter $0$ denotes a $0$ matrix of appropriate size $(n-r) \times (m-r)$. So by using the preceding trick, we may put both $M$ and $M'$ in this form. The two are isomorphic iff $\rk(M) = \rk(M')$, so for PC degree we have now reduced to the case of showing that $I_r \oplus 0$ and $I_{r'} \oplus 0$ are not isomorphic when $r \neq r'$. 

Note that, aside from the equations saying $X$ and $Y$ are invertible, this is almost identical to the Rank Principle (see Section~\ref{sec:prelim:inv}). In the rest of this section we will prove PC lower bounds on both the Rank Principle and the Inversion Principle. Here, we show that the addition of these invertibility axioms in fact makes 2TI much easier in PC than the Rank Principle or 3TI.

\begin{proposition}
Let $M,M'$ be two $n \times m$ matrices of  ranks $r,r'$ respectively, with $r' > r$. Then, over any field whose characteristic does not divide $r'-r$, the following equations have a degree 3 PC refutation and a degree 4 NS refutation:
\[
XMY^T = M' \qquad XX' = X' X = \Id_n \qquad YY' = Y' Y = \Id_m.
\]
\end{proposition}

For those familiar with the low-degree PC proof of the functional onto-PHP, the following proof is similar. 
\newcommand{\Tr}{\text{Tr}}

\begin{proof}[Proof idea]
By the observations in Section~\ref{sec:trick}, we may assume without loss of generality (from the point of view of PC degree) that $M = \Id_r \oplus 0_{n - r \times m-r}$ and $M' = \Id_{r'} \oplus 0_{n-r' \times m-r'}$. 

Write $X = \begin{bmatrix} X_{11} & X_{12} \\ X_{21} & X_{22} \end{bmatrix}$ where the top-left block $X_{11}$ has size $r' \times r$, and similarly write $Y = \begin{bmatrix} Y_{11} & Y_{12} \\ Y_{21} & Y_{22} \end{bmatrix}$ where $Y_{11}$ has size $r' \times r$. In this notation, the matrix equation $XMY^T = M'$ becomes the equations
\begin{align*}
XMY^T & = \begin{bmatrix} X_{11} & X_{12} \\ X_{21} & X_{22} \end{bmatrix} \begin{bmatrix} \Id_r \\ & 0_{(n-r) \times (m-r)} \end{bmatrix} \begin{bmatrix} Y_{11}^T & Y_{21}^T \\ Y_{12}^T & Y_{22}^T  \end{bmatrix} \\
 & = \begin{bmatrix} X_{11} & 0 \\ X_{21} & 0 \end{bmatrix} \begin{bmatrix} Y_{11}^T & Y_{21}^T \\ 0 & 0 \end{bmatrix} \\
 & = M' = \begin{bmatrix} \Id_{r'} \\ & 0_{(n-r') \times (m-r')} \end{bmatrix}
\end{align*}
which becomes the four matrix equations
\begin{equation} \label{eq:2}
X_{11} Y_{11}^T = \Id_{r'} \qquad X_{11}Y_{21}^T = 0 \qquad X_{21}Y_{11}^T = 0 \qquad X_{21}Y_{21}^T = 0.
\end{equation}
Note that so far our PC proof hasn't actually done anything---it is all just notation, and all in the same degree we started with (degree 2).

Then, using the equations $XX' = \Id_n$ and $YY'=\Id_m$, we will derive that $Y_{11}^T X_{11} = \Id_r$. Then we derive 1 as 
\[
\frac{1}{r-r'}\left(\Tr(X_{11} Y_{11}^T - \Id_{r'}) - \Tr(Y_{11}^T X_{11} - \Id_r)\right).
\]
The point here is that trace is additive and cyclically invariant, so $\Tr(X_{11} Y_{11}^T) \equiv \Tr(Y_{11}^T X_{11})$, identically as polynomials, so there is no further derivation needed.
\end{proof}

\begin{proof}
The proof starts using the first part of the proof idea above, so we continue from Equation~(\ref{eq:2}) with the notation introduced above.
In the remainder of the proof, we will derive $Y_{11}^T X_{11} = \Id_r$. Then the last paragraph of the proof idea will complete the proof.

To derive $Y_{11}^T X_{11} = \Id_r$, we will use the invertibility equations (those involving $X'$ and $Y'$). Write $X' = \begin{bmatrix} X'_{11} & X'_{12} \\ X'_{21} & X'_{22} \end{bmatrix}$, where $X'_{11}$ has size $r \times r'$ (NB: the size is the ``transpose'' of the size of $X_{11}$) and similarly for $Y'$. 

From considering the upper-left $r \times r$ block of the matrix equation $X' X = \Id_n$, we get
\[
X'_{11} X_{11} + X'_{12} X_{21} = \Id_r.
\]
Right multiplying by $Y_{11}^T$, we get
\[
X'_{11} X_{11} Y_{11}^T + X'_{12} X_{21} Y_{11}^T = Y_{11}^T. 
\]
But now we can subtract from this $X'_{11}$ times the equation $X_{11} Y_{11}^T = \Id_{r'}$, and also $X'_{12}$ times the equation $X_{21} Y_{11}^T = 0$ to get
\begin{equation}
X'_{11} = Y_{11}^T.
\end{equation}

Similarly, considering the upper-left $r \times r$ block of the matrix equation $Y' Y = \Id_m$, we get $Y'_{11} Y_{11} + Y'_{12} Y_{21} = \Id_{r}$. For consistency with the notation above, we take the transpose of this entire equation (in PC, this is essentially a null-op---we are just re-arranging how we are viewing a set of $(r')^2$ equations on the page), to get:
\[
Y_{11}^T (Y'_{11})^T + Y_{21}^T (Y'_{12})^T = \Id_r.
\]
Left multiplying by $X_{11}$, we get
\[
X_{11} Y_{11}^T (Y'_{11})^T + X_{11} Y_{21}^T (Y'_{12})^T = X_{11}.
\]
Now, right-multiplying the equation $X_{11} Y_{11}^T = \Id$ by $(Y'_{11})^T$, and right-multiplying the equation $X_{11} Y_{21}^T = 0$ by $(Y'_{12})^T$ and subtracting both of these from the above, we get
\begin{equation}
(Y'_{11})^T = X_{11}.
\end{equation}

Next, we derive $M - X' M' (Y')^T=0$, as follows: left-multiplying $XMY^T - M'$ by $X'$, and subtract from it $(X'X-I)$ times $MY^T$, to get $-X'M' + MY^T$. Now right-multiply the latter by $(Y')^T$ and subtract from it $M$ times $(Y^T (Y')^T - I)$, yielding $-X'M'(Y')^T + M$. Now multiply by $-1$.

Now, from $X' M' (Y')^T = M$, as at the beginning of the proof, we derive that $X'_{11} (Y'_{11})^T = \Id_r$. But above we have derived that $X'_{11} = Y_{11}^T$ and $(Y'_{11})^T = X_{11}$, so from the preceding three equations we get $Y_{11}^T X_{11} = \Id_r$, as claimed. This completes the PC proof.

Let us unroll the PC proof to derive a Nullstellensatz proof (here we underline the use of original equations):
\begin{align*}
r - r' & = \Tr(\underline{X_{11} Y_{11}^T - \Id_{r'}}) - \Tr(Y_{11}^T X_{11} - \Id_r)  \\
\end{align*}
Now we focus on the NS derivation of $Y_{11}^T X_{11} - \Id_r$. Since the trace is linear, and we are focusing on degree, this is without loss of generality. We have:
\begin{align*}
Y_{11}^T X_{11} - \Id_r = & \left(X'_{11} (Y'_{11})^T - \Id_r\right) - (X'_{11} - Y_{11}^T) (Y'_{11})^T -  Y_{11}^T \left((Y'_{11})^T - X_{11}\right) \\
 = & \left(-X' \underline{(XMY^T - M')} Y'^T + \underline{(X'X - I)}MY^T (Y')^T + M \underline{(Y^T (Y')^T - I)}\right)_{11} \\
  & + \left(\underline{(X'_{11} X_{11} + X'_{12} X_{21} - \Id_r)} Y_{11}^T - X'_{11} \underline{(X_{11} Y_{11}^T - \Id_{r'})} - X'_{12}\underline{(X_{21} Y_{11}^T)} \right) (Y'_{11})^T \\ 
  & -  Y_{11}^T \left( X_{11} \underline{(Y_{11}^T (Y'_{11})^T + Y_{21}^T (Y'_{12})^T - \Id_r)} - \underline{(X_{11} Y_{11}^T - \Id)} (Y'_{11})^T - \underline{(X_{11} Y_{21}^T)} (Y'_{12})^T   \right). 
\end{align*}
This is visibly degree 4.
\end{proof}

\subsection{Inversion Principle implies the Rank Principle}

\begin{lemma} \label{lem:inv_to_rank}
If the $r \times r$ Inversion Principle has a degree $d$ PC derivation, 
then there is a degree $\max\{d,3\}$ PC refutation of the Rank Principle stating that a rank $r$ matrix is not equivalent (isomorphic) to a rank $n$ matrix, for any $n > r$. 

If the Inversion Principle has a degree $d$ NS derivation, then the Rank Principle has a degree $d+2$ NS refutation.
\end{lemma}

\begin{proof}
Suppose the $r \times r$ Inversion Principle has a degree-$d$ derivation. Consider the Rank Principle $XY=I_n$ where $X$ is $n \times r$ and $Y$ is $r \times n$, with $n > r$. Write
\[
X = \begin{bmatrix} X_0 \\ X_1 \end{bmatrix} \text{ and } Y = \begin{bmatrix} Y_0 & Y_1 \end{bmatrix},
\]
where $X_0,Y_0$ are $r \times r$. Then, examining the upper-left $r \times r$ corner of the original equations, we find $X_0 Y_0 = I_r$. As these are square matrices, by assumption in degree $d$ we may then derive that $Y_0 X_0 = I_r$ as well.

Now, multiply both sides of $XY=I_n$ on the left by the matrix $\begin{bmatrix} Y_0 & 0 \\ 0 & I_{n-r} \end{bmatrix}$. The result is then the set of degree-3 equations
\[
\begin{bmatrix} Y_0 X_0 \\ X_1 \end{bmatrix} \begin{bmatrix} Y_0 & Y_1 \end{bmatrix} = \begin{bmatrix} Y_0 & 0 \\ 0 & I_{n-r} \end{bmatrix}.
\]
Considering the upper-right $r \times (n-r)$ block of these equations, we find the equations $Y_0 X_0 Y_1 = 0$. 

But now, from the equation $Y_0 X_0 = I_r$, we may right-multiply by $Y_1$ to get $Y_0 X_0 Y_1 = Y_1$. Combining with the equation at the end of the last paragraph, we then conclude $Y_1 = 0$.

Finally, consider the lower-right $(n-r) \times (n-r)$ part of the original equation $XY = I_n$, namely, $X_1 Y_1 = I_{n-r}$. We had already derived $Y_1 = 0$, which we can then left-multiply by $X_1$ to get $X_1 Y_1 = 0$. Considering any diagonal entry of these two equations, we then derive the contradiction $1=0$.

To see the NS certificate, we unwrap the above proof. First write $Y_0 X_0 - I_r$ as a linear combination of the equations $X_0 Y_0 - I_r$ with polynomial coefficients, in total degree $d$. Among our starting equations in the Rank Principle, we have $X_0 Y_1$ and $X_1 Y_1 - I_{n-r}$. Then the following linear combination has degree 2 more than $Y_0 X_0 - I_r$, and derives $1$ in any of its diagonal entries: 
\[
-X_1 Y_0 \underline{X_0 Y_1} + X_1 (Y_0 X_0 - I_r) Y_1 + \underline{(X_1 Y_1 - I_{n-r})}.
\]
\end{proof}

\begin{observation}
The $n \times n$ Inversion Principle has a proof of degree $2n+2$.
\end{observation}

\begin{proof}
The idea is to use Laplace expansion. We spell out the details.

We start with $XY = I_n$, where $X$ and $Y$ are $n \times n$ matrices of variables. Left-multiply by $Y$ to get $YXY = Y$, and then right multiply by $Adj(Y)$ (whose entries are the $(n-1) \times (n-1)$ cofactors of $Y$, hence have degree $n-1$) to get $YXYAdj(Y) = YAdj(Y)$. Now, by Laplace expansion, we have $Y Adj(Y) \equiv \det(Y) I_n$, so we get $YX \det(Y) = \det(Y) I_n$.

Next, starting from $XY=I_n$ and expanding out the determinant term-by-term, we derive $\det(XY) = 1$. (Note that here, we are not simply applying the determinant to the matrix $XY-I$, as that would give us the value of the characteristic polynomial evaluated at 1. Instead, we repeatedly use that from $a-b=0$ and $c-d=0$ we can derive $ac-bd=0$ as $\underline{(a-b)}c + b\underline{(c-d)}$. Similarly, we can derive $(a+c)-(b+d)=0$ as $(a-b) + (c-d)$.) Now, since $\det(XY) \equiv \det(X)\det(Y)$ identically as polynomials, we have derived $\det(X)\det(Y)=1$ in degree $n$.

Now, from $YX \det(Y) - \det(Y) I_n$ in the first paragraph, we multiply by $\det(X)$ to get $(YX - I_n)(\det(X)\det(Y))$. From $\det(X)\det(Y)-1$ in the second paragraph, we multiply by $-(YX-I_n)$ and add to the preceding to get $YX - I_n$, all in degree at most $2n+2$.
\end{proof}


\subsection{Lower bound on the Rank Principle (and Inversion Principle) via reduction from PHP}
\newcommand{\bbI}{\mathbb{I}}

Here we show that the Rank Principle (see Section~\ref{sec:prelim:inv}) requires large PC degree, via a reduction to the Pigeonhole Principle. 
For the Pigeonhole principle, a tight PC degree lower bound is known:


\begin{theorem}[{Razborov \cite{DBLP:journals/cc/Razborov98}}]
\label{thm:PCdeglb}
Any $PC$ refutation of the Functional $PHP^n_r$ requires degree $r/2$+1 over any field.
\end{theorem}

We use this to show:

\begin{theorem} \label{thm:rank}
Let $n\in \mathbb N$,  $n\geq 2$ and $1\leq r<n$. $\bbI(r,n)$ (with or without the Boolean axioms) requires degree $r/2+1$ in $PC$ over any field.
\end{theorem}
\begin{proof}
We prove that $PHP^n_{r}$ is $(1,2)$-reducible to $\bbI(r,n)$. 
First we consider the following degree $1$  polynomials  defining $x$ and $y$ variables of $\bbI(r,n)$ 
in terms of the $p$ variables of $PHP^n_{r}$.
variables 
$$
x_{i,k} = y_{i,k} = p_{i,k} \qquad \mbox{ for } i \in [n], k \in [n-1].
$$ 
Second we show a degree $2$ $PC$ proof   of $\bbI(r,n)$ from the  polynomials defining the  $PHP^n_r$.
From $PHP$ axioms $p_{i,k}p_{k,j}$ for $i,j \in [n], i\not =j$,  and summing over all $k \in [r]$, we get 
$$
\sum_{k \in [r]}p_{i,k}p_{k,j},
$$ 
which are exactly the axioms of $\bbI(r,n)$ for $i \not =j, i,j \in [n]$, after the substitution of variables.

For a $i\in [n]$, take  the boolean axioms written in the form $p_{i,k}p_{i,k}-p_{i,k}$ and  sum them over $k \in [r]$:
$$\sum_{k \in [r]} p_{i,k}p_{i,k} - \sum_{k \in [r]}p_{i,k}$$
Summing this last polynomial with the $PHP$ axiom $\sum_{k \in [r]} p_{i,k} -1$
we get the polynomial
$$\sum_{k\in [r]}p_{i,k}p_{i,k} -1,$$
which is the axiom of $\bbI(r,n)$ for $i=j$ after the substitution of the  variables.  The proof has degree $2$. 
The result follows immediately  from Lemma \ref{lem:PCred} and Theorem \ref{thm:PCdeglb}. 

\end{proof}

\begin{corollary}
Any $PC$ proof of $AB=I \vdash BA=I$, where $A,B$ are square $n\times n$ $\{0,1\}$ matrices requires degree $n/2+1$.
\end{corollary}

\begin{proof}
Follows immediately from Theorem~\ref{thm:rank} and Lemma~\ref{lem:inv_to_rank}.
\end{proof}

\newcommand{\etalchar}[1]{$^{#1}$}
\begin{thebibliography}{OWWZ14}

\bibitem[ABRW04]{ABRW}
Michael Alekhnovich, Eli Ben{-}Sasson, Alexander~A. Razborov, and Avi
  Wigderson.
\newblock Pseudorandom generators in propositional proof complexity.
\newblock {\em {SIAM} J. Comput.}, 34(1):67--88, 2004.
\newblock \href {https://doi.org/10.1137/S0097539701389944}
  {\path{doi:10.1137/S0097539701389944}}.

\bibitem[AM13]{AM}
Albert Atserias and Elitza~N. Maneva.
\newblock {Sherali}--{Adams} relaxations and indistinguishability in counting
  logics.
\newblock {\em {SIAM} J. Comput.}, 42(1):112--137, 2013.
\newblock \href {https://doi.org/10.1137/120867834}
  {\path{doi:10.1137/120867834}}.

\bibitem[AS05]{AS05}
Manindra Agrawal and Nitin Saxena.
\newblock Automorphisms of finite rings and applications to complexity of
  problems.
\newblock In {\em {STACS} 2005, 22nd Annual Symposium on Theoretical Aspects of
  Computer Science, Proceedings}, pages 1--17, 2005.
\newblock \href {https://doi.org/10.1007/978-3-540-31856-9_1}
  {\path{doi:10.1007/978-3-540-31856-9_1}}.

\bibitem[Bab16]{Babai}
L\'{a}szl\'{o} Babai.
\newblock Graph isomorphism in quasipolynomial time [extended abstract].
\newblock In {\em S{TOC}'16---{P}roceedings of the 48th {A}nnual {ACM} {SIGACT}
  {S}ymposium on {T}heory of {C}omputing}, pages 684--697. ACM, New York, 2016.
\newblock \href {https://doi.org/10.1145/2897518.2897542}
  {\path{doi:10.1145/2897518.2897542}}.

\bibitem[BCIP02]{BCIP}
Josh Buresh{-}Oppenheim, Matthew Clegg, Russell Impagliazzo, and Toniann
  Pitassi.
\newblock Homogenization and the polynomial calculus.
\newblock {\em Comput. Complex.}, 11(3-4):91--108, 2002.
\newblock \href {https://doi.org/10.1007/s00037-002-0171-6}
  {\path{doi:10.1007/s00037-002-0171-6}}.

\bibitem[Ber18]{berkholz2018relation}
Christoph Berkholz.
\newblock The relation between polynomial calculus, sherali-adams, and
  sum-of-squares proofs.
\newblock In {\em 35th Symposium on Theoretical Aspects of Computer Science
  (STACS 2018)}. Schloss Dagstuhl-Leibniz-Zentrum fuer Informatik, 2018.

\bibitem[BG15]{BerkholzG15}
Christoph Berkholz and Martin Grohe.
\newblock Limitations of algebraic approaches to graph isomorphism testing.
\newblock In Magn{\'{u}}s~M. Halld{\'{o}}rsson, Kazuo Iwama, Naoki Kobayashi,
  and Bettina Speckmann, editors, {\em Automata, Languages, and Programming -
  42nd International Colloquium, {ICALP} 2015, Kyoto, Japan, July 6-10, 2015,
  Proceedings, Part {I}}, volume 9134 of {\em Lecture Notes in Computer
  Science}, pages 155--166. Springer, 2015.

\bibitem[BG17]{BerkholzG16}
Christoph Berkholz and Martin Grohe.
\newblock Linear diophantine equations, group csps, and graph isomorphism.
\newblock In Philip~N. Klein, editor, {\em Proceedings of the Twenty-Eighth
  Annual {ACM-SIAM} Symposium on Discrete Algorithms, {SODA} 2017, Barcelona,
  Spain, Hotel Porta Fira, January 16-19}, pages 327--339. {SIAM}, 2017.
\newblock Preprint \href{http://arxiv.org/abs/1607.04287}{arXiv:1607.04287
  [cs.CC]}.
\newblock \href {https://doi.org/10.1137/1.9781611974782.21}
  {\path{doi:10.1137/1.9781611974782.21}}.

\bibitem[BGIP01]{DBLP:journals/jcss/BussGIP01}
Samuel~R. Buss, Dima Grigoriev, Russell Impagliazzo, and Toniann Pitassi.
\newblock Linear gaps between degrees for the polynomial calculus modulo
  distinct primes.
\newblock {\em J. Comput. Syst. Sci.}, 62(2):267--289, 2001.
\newblock \href {https://doi.org/10.1006/jcss.2000.1726}
  {\path{doi:10.1006/jcss.2000.1726}}.

\bibitem[BGL{\etalchar{+}}19]{BGLQW}
Peter~A. Brooksbank, Joshua~A. Grochow, Yinan Li, Youming Qiao, and James~B.
  Wilson.
\newblock Incorporating {Weisfeiler}--{Leman} into algorithms for group
  isomorphism.
\newblock \href{https://arxiv.org/abs/1905.02518}{arXiv:1905.02518 [cs.CC]},
  2019.

\bibitem[BI99]{BI}
Eli Ben{-}Sasson and Russell Impagliazzo.
\newblock Random {CNF}'s are hard for the {Polynomial} {Calculus}.
\newblock In {\em 40th Annual Symposium on Foundations of Computer Science,
  {FOCS} '99, 17-18 October, 1999, New York, NY, {USA}}, pages 415--421. {IEEE}
  Computer Society, 1999.
\newblock (Journal version in \emph{Comput. Complex.} 2010,
  doi:10.1007/s00037-010-0293-1).
\newblock \href {https://doi.org/10.1109/SFFCS.1999.814613}
  {\path{doi:10.1109/SFFCS.1999.814613}}.

\bibitem[Bre70]{brent}
R.~P. Brent.
\newblock Algorithms for matrix multiplication.
\newblock Stanford Computer Science Dept. Tech. Report STAN-CS-70-157,
  available online at
  \href{https://apps.dtic.mil/sti/pdfs/AD0705509.pdf}{https://apps.dtic.mil/sti/pdfs/AD0705509.pdf},
  1970.

\bibitem[BS20]{BS1}
Jendrik Brachter and Pascal Schweitzer.
\newblock On the {Weisfeiler}--{Leman} dimension of finite groups.
\newblock In Holger Hermanns, Lijun Zhang, Naoki Kobayashi, and Dale Miller,
  editors, {\em {LICS} '20: 35th Annual {ACM/IEEE} Symposium on Logic in
  Computer Science, Saarbr{\"{u}}cken, Germany, July 8-11, 2020}, pages
  287--300. {ACM}, 2020.
\newblock \href {https://doi.org/10.1145/3373718.3394786}
  {\path{doi:10.1145/3373718.3394786}}.

\bibitem[BS22]{BS2}
Jendrik Brachter and Pascal Schweitzer.
\newblock A systematic study of isomorphism invariants of finite groups via the
  {Weisfeiler}--{Leman} dimension.
\newblock In Shiri Chechik, Gonzalo Navarro, Eva Rotenberg, and Grzegorz
  Herman, editors, {\em 30th Annual European Symposium on Algorithms, {ESA}
  2022, September 5-9, 2022, Berlin/Potsdam, Germany}, volume 244 of {\em
  LIPIcs}, pages 27:1--27:14. Schloss Dagstuhl - Leibniz-Zentrum f{\"{u}}r
  Informatik, 2022.
\newblock \href {https://doi.org/10.4230/LIPIcs.ESA.2022.27}
  {\path{doi:10.4230/LIPIcs.ESA.2022.27}}.

\bibitem[CEI96]{CEI}
Matthew Clegg, Jeffery Edmonds, and Russell Impagliazzo.
\newblock Using the {G}roebner basis algorithm to find proofs of
  unsatisfiability.
\newblock In {\em Proceedings of the {T}wenty-eighth {A}nnual {ACM} {S}ymposium
  on the {T}heory of {C}omputing ({P}hiladelphia, {PA}, 1996)}, pages 174--183.
  ACM, New York, 1996.
\newblock \href {https://doi.org/10.1145/237814.237860}
  {\path{doi:10.1145/237814.237860}}.

\bibitem[CFI92]{CFI}
Jin-Yi Cai, Martin F\"{u}rer, and Neil Immerman.
\newblock An optimal lower bound on the number of variables for graph
  identification.
\newblock {\em Combinatorica}, 12(4):389--410, 1992.
\newblock \href {https://doi.org/10.1007/BF01305232}
  {\path{doi:10.1007/BF01305232}}.

\bibitem[CL22]{CL}
Nathaniel~A. Collins and Michael Levet.
\newblock Count-free {Weisfeiler}--{Leman} and group isomorphism.
\newblock \href{https://arxiv.org/abs/2212.11247}{arXiv:2212.11247 [cs.DS]},
  2022.

\bibitem[FGS19]{FGS}
Vyacheslav Futorny, Joshua~A. Grochow, and Vladimir~V. Sergeichuk.
\newblock Wildness for tensors.
\newblock {\em Linear Algebra Appl.}, 566:212--244, 2019.
\newblock \href {https://doi.org/10.1016/j.laa.2018.12.022}
  {\path{doi:10.1016/j.laa.2018.12.022}}.

\bibitem[FP06]{DBLP:conf/eurocrypt/FaugereP06}
Jean{-}Charles Faug{\`{e}}re and Ludovic Perret.
\newblock Polynomial equivalence problems: Algorithmic and theoretical aspects.
\newblock In Serge Vaudenay, editor, {\em Advances in Cryptology - {EUROCRYPT}
  2006, 25th Annual International Conference on the Theory and Applications of
  Cryptographic Techniques, St. Petersburg, Russia, May 28 - June 1, 2006,
  Proceedings}, volume 4004 of {\em Lecture Notes in Computer Science}, pages
  30--47. Springer, 2006.
\newblock \href {https://doi.org/10.1007/11761679\_3}
  {\path{doi:10.1007/11761679\_3}}.

\bibitem[GQ21a]{GQ2}
Joshua~A. Grochow and Youming Qiao.
\newblock On p-group isomorphism: Search-to-decision, counting-to-decision, and
  nilpotency class reductions via tensors.
\newblock In Valentine Kabanets, editor, {\em 36th Computational Complexity
  Conference, {CCC} 2021, July 20-23, 2021, Toronto, Ontario, Canada (Virtual
  Conference)}, volume 200 of {\em LIPIcs}, pages 16:1--16:38. Schloss Dagstuhl
  - Leibniz-Zentrum f{\"{u}}r Informatik, 2021.
\newblock \href {https://doi.org/10.4230/LIPIcs.CCC.2021.16}
  {\path{doi:10.4230/LIPIcs.CCC.2021.16}}.

\bibitem[GQ21b]{GQ}
Joshua~A. Grochow and Youming Qiao.
\newblock On the complexity of isomorphism problems for tensors, groups, and
  polynomials {I:} tensor isomorphism-completeness.
\newblock In James~R. Lee, editor, {\em 12th Innovations in Theoretical
  Computer Science Conference, {ITCS} 2021, January 6-8, 2021, Virtual
  Conference}, volume 185 of {\em LIPIcs}, pages 31:1--31:19. Schloss Dagstuhl
  - Leibniz-Zentrum f{\"{u}}r Informatik, 2021.
\newblock \href {https://doi.org/10.4230/LIPIcs.ITCS.2021.31}
  {\path{doi:10.4230/LIPIcs.ITCS.2021.31}}.

\bibitem[Gri81]{grigorievGI}
D.~Ju. Grigoriev.
\newblock Complexity of ``wild'' matrix problems and of the isomorphism of
  algebras and graphs.
\newblock {\em Zap. Nauchn. Sem. Leningrad. Otdel. Mat. Inst. Steklov. (LOMI)},
  105:10--17, 198, 1981.
\newblock Theoretical applications of the methods of mathematical logic, III.
\newblock \href {https://doi.org/10.1007/BF01084390}
  {\path{doi:10.1007/BF01084390}}.

\bibitem[Gri01]{grigoriev2001linear}
Dima Grigoriev.
\newblock Linear lower bound on degrees of positivstellensatz calculus proofs
  for the parity.
\newblock {\em Theoretical Computer Science}, 259(1-2):613--622, 2001.

\bibitem[Gri13]{grigoriev}
Dima Grigoriev.
\newblock Polynomial complexity of solving systems of few algebraic equations
  with small degrees.
\newblock In Vladimir~P. Gerdt, Wolfram Koepf, Ernst~W. Mayr, and Evgenii~V.
  Vorozhtsov, editors, {\em Computer Algebra in Scientific Computing - 15th
  International Workshop, {CASC} 2013, Berlin, Germany, September 9-13, 2013.
  Proceedings}, volume 8136 of {\em Lecture Notes in Computer Science}, pages
  136--139. Springer, 2013.
\newblock \href {https://doi.org/10.1007/978-3-319-02297-0\_11}
  {\path{doi:10.1007/978-3-319-02297-0\_11}}.

\bibitem[Gro12]{GrochowLie}
Joshua~A. Grochow.
\newblock Matrix {Lie} algebra isomorphism.
\newblock In {\em IEEE Conference on Computational Complexity (CCC12)}, pages
  203--213, 2012.
\newblock Also available as arXiv:1112.2012 [cs.CC] and ECCC Technical Report
  TR11-168.
\newblock \href {https://doi.org/10.1109/CCC.2012.34}
  {\path{doi:10.1109/CCC.2012.34}}.

\bibitem[Gro13]{GrochowSE}
Joshua~A. Grochow.
\newblock Answer to ``deciding bound on tensor rank for a fixed value''.
\newblock CSTheory StackExchange,
  \href{https://cstheory.stackexchange.com/a/19518/129}{https://cstheory.stackexchange.com/a/19518/129},
  2013.

\bibitem[Hel19]{helfgott}
Harald~Andr\'{e}s Helfgott.
\newblock Isomorphismes de graphes en temps quasi-polynomial [d'apr\`es {B}abai
  et {L}uks, {W}eisfeiler-{L}eman,\dots ].
\newblock {\em Ast\'{e}risque}, (407):Exp. No. 1125, 135--182, 2019.
\newblock S\'{e}minaire Bourbaki. Vol. 2016/2017. Expos\'{e}s 1120--1135.
  English translation with appendices by Jitendra Bajpai and Daniele Dona
  available at \href{https://arxiv.org/abs/1710.04574}{arXiv:17010.04574
  [math.GR]}.
\newblock \href {https://doi.org/10.24033/ast} {\path{doi:10.24033/ast}}.

\bibitem[HQ21]{HQ}
Xiaoyu He and Youming Qiao.
\newblock On the {B}aer-{L}ov\'{a}sz-{T}utte construction of groups from
  graphs: isomorphism types and homomorphism notions.
\newblock {\em European J. Combin.}, 98:Paper No. 103404, 12, 2021.
\newblock \href {https://doi.org/10.1016/j.ejc.2021.103404}
  {\path{doi:10.1016/j.ejc.2021.103404}}.

\bibitem[HT15]{HT}
Pavel Hrubes and Iddo Tzameret.
\newblock Short proofs for the determinant identities.
\newblock {\em {SIAM} J. Comput.}, 44(2):340--383, 2015.
\newblock \href {https://doi.org/10.1137/130917788}
  {\path{doi:10.1137/130917788}}.

\bibitem[JQSY19]{JQSY19}
Zhengfeng Ji, Youming Qiao, Fang Song, and Aaram Yun.
\newblock General linear group action on tensors: {A} candidate for
  post-quantum cryptography.
\newblock In Dennis Hofheinz and Alon Rosen, editors, {\em Theory of
  Cryptography - 17th International Conference, {TCC} 2019, Nuremberg, Germany,
  December 1-5, 2019, Proceedings, Part {I}}, volume 11891 of {\em Lecture
  Notes in Computer Science}, pages 251--281. Springer, 2019.
\newblock Preprint \href{https://arxiv.org/abs/1906.04330}{arXiv:1906.04330
  [cs.CR]}.
\newblock \href {https://doi.org/10.1007/978-3-030-36030-6\_11}
  {\path{doi:10.1007/978-3-030-36030-6\_11}}.

\bibitem[Lan12]{landsbergBook}
J.~M. Landsberg.
\newblock {\em Tensors: geometry and applications}, volume 128 of {\em Graduate
  Studies in Mathematics}.
\newblock American Mathematical Society, Providence, RI, 2012.
\newblock \href {https://doi.org/10.1090/gsm/128} {\path{doi:10.1090/gsm/128}}.

\bibitem[Las01]{lasserre}
Jean~B. Lasserre.
\newblock Global optimization with polynomials and the problem of moments.
\newblock {\em SIAM J. Optim.}, 11(3):796--817, 2000/01.
\newblock \href {https://doi.org/10.1137/S1052623400366802}
  {\path{doi:10.1137/S1052623400366802}}.

\bibitem[Lic85]{lickteigTypical}
Thomas Lickteig.
\newblock Typical tensorial rank.
\newblock {\em Linear Algebra Appl.}, 69:95--120, 1985.
\newblock \href {https://doi.org/10.1016/0024-3795(85)90070-9}
  {\path{doi:10.1016/0024-3795(85)90070-9}}.

\bibitem[LL89]{LL}
Thomas Lehmkuhl and Thomas Lickteig.
\newblock On the order of approximation in approximative triadic decompositions
  of tensors.
\newblock {\em Theoret. Comput. Sci.}, 66(1):1--14, 1989.
\newblock \href {https://doi.org/10.1016/0304-3975(89)90141-2}
  {\path{doi:10.1016/0304-3975(89)90141-2}}.

\bibitem[Luk93]{Luks}
Eugene~M. Luks.
\newblock Permutation groups and polynomial-time computation.
\newblock In {\em Groups and computation ({N}ew {B}runswick, {NJ}, 1991)},
  volume~11 of {\em DIMACS Ser. Discrete Math. Theoret. Comput. Sci.}, pages
  139--175. Amer. Math. Soc., Providence, RI, 1993.

\bibitem[McK81]{mckay}
Brendan~D. McKay.
\newblock Practical graph isomorphism.
\newblock {\em Congr. Numer.}, 30:45--87, 1981.

\bibitem[Miy96]{Miyazaki}
Takunari Miyazaki.
\newblock {Luks's} reduction of graph isomorphism to code equivalence.
\newblock Comment to E. W. Clark,
  \href{https://groups.google.com/forum/#!msg/sci.math.research/puZxGj9HXKI/CeyH2yyyNFUJ}{https://groups.google.com/forum/\#!msg/sci.math.research/puZxGj9HXKI/CeyH2yyyNFUJ},
  1996.

\bibitem[MP14]{MP}
Brendan~D. McKay and Adolfo Piperno.
\newblock Practical graph isomorphism, {II}.
\newblock {\em J. Symbolic Comput.}, 60:94--112, 2014.
\newblock \href {https://doi.org/10.1016/j.jsc.2013.09.003}
  {\path{doi:10.1016/j.jsc.2013.09.003}}.

\bibitem[NPOV15]{TNN}
Alexander Novikov, Dmitry Podoprikhin, Anton Osokin, and Dmitry Vetrov.
\newblock Tensorizing neural networks.
\newblock In {\em Proceedings of the 28th International Conference on Neural
  Information Processing Systems - Volume 1}, NIPS'15, pages 442--450. MIT
  Press, 2015.

\bibitem[OWWZ14]{DBLP:conf/soda/ODonnellWWZ14}
Ryan O'Donnell, John Wright, Chenggang Wu, and Yuan Zhou.
\newblock Hardness of robust graph isomorphism, lasserre gaps, and asymmetry of
  random graphs.
\newblock In Chandra Chekuri, editor, {\em Proceedings of the Twenty-Fifth
  Annual {ACM-SIAM} Symposium on Discrete Algorithms, {SODA} 2014, Portland,
  Oregon, USA, January 5-7, 2014}, pages 1659--1677. {SIAM}, 2014.
\newblock Preprint available as
  arXiv:\href{https://arxiv.org/abs/1401.2436}{1401.2436 [cs.CC]}.
\newblock \href {https://doi.org/10.1137/1.9781611973402.120}
  {\path{doi:10.1137/1.9781611973402.120}}.

\bibitem[Pat96]{Pat96}
Jacques Patarin.
\newblock Hidden fields equations {(HFE)} and isomorphisms of polynomials
  {(IP):} two new families of asymmetric algorithms.
\newblock In {\em Advances in Cryptology - {EUROCRYPT} '96, International
  Conference on the Theory and Application of Cryptographic Techniques,
  Saragossa, Spain, May 12-16, 1996, Proceeding}, pages 33--48, 1996.
\newblock \href {https://doi.org/10.1007/3-540-68339-9_4}
  {\path{doi:10.1007/3-540-68339-9_4}}.

\bibitem[PR97]{PR}
Erez Petrank and Ron~M. Roth.
\newblock Is code equivalence easy to decide?
\newblock {\em {IEEE} Trans. Inf. Theory}, 43(5):1602--1604, 1997.
\newblock \href {https://doi.org/10.1109/18.623157}
  {\path{doi:10.1109/18.623157}}.

\bibitem[Raz98]{DBLP:journals/cc/Razborov98}
Alexander~A. Razborov.
\newblock Lower bounds for the polynomial calculus.
\newblock {\em Comput. Complex.}, 7(4):291--324, 1998.
\newblock \href {https://doi.org/10.1007/s000370050013}
  {\path{doi:10.1007/s000370050013}}.

\bibitem[SC04]{SC}
Michael Soltys and Stephen Cook.
\newblock The proof complexity of linear algebra.
\newblock {\em Ann. Pure Appl. Logic}, 130(1-3):277--323, 2004.
\newblock \href {https://doi.org/10.1016/j.apal.2003.10.018}
  {\path{doi:10.1016/j.apal.2003.10.018}}.

\bibitem[Sch08]{schoenebeck2008linear}
Grant Schoenebeck.
\newblock Linear level lasserre lower bounds for certain k-csps.
\newblock In {\em 2008 49th Annual IEEE Symposium on Foundations of Computer
  Science}, pages 593--602. IEEE, 2008.

\bibitem[Sol01]{soltys}
Michael Soltys.
\newblock {\em The complexity of derivations of matrix identities}.
\newblock PhD thesis, University of Toronto, 2001.
\newblock Availalble on ECCC at
  \href{https://eccc.weizmann.ac.il/resources/pdf/soltys.pdf}{https://eccc.weizmann.ac.il/resources/pdf/soltys.pdf}.

\bibitem[Som99]{sombra}
Mart\'{\i}n Sombra.
\newblock A sparse effective {N}ullstellensatz.
\newblock {\em Adv. in Appl. Math.}, 22(2):271--295, 1999.
\newblock \href {https://doi.org/10.1006/aama.1998.0633}
  {\path{doi:10.1006/aama.1998.0633}}.

\bibitem[SSC14]{SSC}
Aaron Snook, Grant Schoenebeck, and Paolo Codenotti.
\newblock Graph {Isomorphism} and the {Lasserre} hierarchy.
\newblock \href{https://arxiv.org/abs/1401.0758}{arXiv:1401.0758 [cs.CC]},
  2014.

\bibitem[TDJ{\etalchar{+}}22]{DBLP:conf/eurocrypt/TangDJPQS22}
Gang Tang, Dung~Hoang Duong, Antoine Joux, Thomas Plantard, Youming Qiao, and
  Willy Susilo.
\newblock Practical post-quantum signature schemes from isomorphism problems of
  trilinear forms.
\newblock In Orr Dunkelman and Stefan Dziembowski, editors, {\em Advances in
  Cryptology - {EUROCRYPT} 2022 - 41st Annual International Conference on the
  Theory and Applications of Cryptographic Techniques, Trondheim, Norway, May
  30 - June 3, 2022, Proceedings, Part {III}}, volume 13277 of {\em Lecture
  Notes in Computer Science}, pages 582--612. Springer, 2022.
\newblock \href {https://doi.org/10.1007/978-3-031-07082-2\_21}
  {\path{doi:10.1007/978-3-031-07082-2\_21}}.

\bibitem[Zui17]{zuiddam}
Jeroen Zuiddam.
\newblock A note on the gap between rank and border rank.
\newblock {\em Linear Algebra Appl.}, 525:33--44, 2017.
\newblock \href {https://doi.org/10.1016/j.laa.2017.03.015}
  {\path{doi:10.1016/j.laa.2017.03.015}}.

\end{thebibliography}
\end{document}
